\chapter{Resultados}\label{resultados}

Los resultados del presente trabajo deben ser abordados desde diferentes puntos de vista, por un lado tenemos los referidos al diseño del producto y las decisiones que se tomaron para ello, por otro lado los aspectos relacionados específicamente a inteligencia artificial y por último la gestión del producto como un todo.

El producto obtenido para el Ministerio Público Fiscal ya cuenta con las partes necesarias para ser utilizado por los usuarios según como se definió inicialmente. Como se definió en los objetivos, se alcanzó un nivel de madurez tecnológica equivalente a TRL3, es decir que si bien el producto no fue probado en un entorno real operativo, se validaron todas sus partes y componentes fundamentales.

Respecto de los modelos utilizados para generar el \eng{forecasting} se puede resumir los resultados de la siguiente manera: Se tiene un modelo que consume todos los datos provistos por el MPF, la selección del mismo fue por múltiples criterios, como por ejemplo la velocidad de cómputo o la escalabilidad. Uno de los criterios más importantes para definirlo fue la explicabilidad del mismo, los árboles de decisión permiten comprender cómo es que el modelo está generando las predicciones. Respecto al rendimiento se alcanzó una performance adecuada para constituir una línea de base, es decir, permite obtener un nivel de exactitud útil y aceptable a partir del cual se podrá escalar mejorando los modelos y agregando nuevos datos. Además, el presente trabajo fue motivo de publicaciones, una para la Escuela de Perfeccionamiento en Investigación Operativa (2026), y otra para las jornadas de SINERGIA (2026), ambas en proceso de revisión a la fecha de elaboración de este documento.

A nivel gestión del producto y del proyecto se pueden destacar varias dimensiones que se fueron abordando a medida que este crecía. Se trabajó con un producto de punta a punta, es decir, comenzando por los requerimientos hasta obtener todas las partes de un sistema funcional. Para ello se investigó el qué y el por qué de cada proceso, de cada persona y de cada tecnología, para luego estudiar cómo la interacción de estos aspectos generan restricciones que deben ser superadas. El resultado final de este proceso es un producto que permite que un usuario no técnico disponga de funcionalidades impulsadas por inteligencia artificial que tendrán un impacto directo en la toma de decisiones en su día a día.
